% !TeX spellcheck = en_US
\chapter{Conclusion and Future Work}\label{ch:conclusion}%

\section{Conclusion}\label{sec:conclusion}

This thesis presented a framework for learning long-horizon cyclic nut threading with a dexterous hand--arm system. Rather than learning the entire task with a monolithic policy, the method decomposes demonstration generation into a reinforcement-learning turning skill, learned in a simplified thread simulation with a privileged asymmetric actor--critic, and scripted release-and-return transitions, and distills the resulting long-horizon trajectories into a unified visuomotor policy for deployment. Human corrections during real-world replay adapt the demonstration distribution to the contact and perception discrepancies that the simulation does not capture, at the cost of overwriting about six percent of the replayed frames.

The student is a conditional flow-matching policy. The thesis derived its objective from the conditional flow-matching construction, showed how it applies to action chunks conditioned on a point-cloud feature, proprioception, a global progress scalar, and a process-phase encoding, and documented the trajectory-Transformer velocity field, the Euler sampler, the parameter-space and action-space averaging, and the Cartesian action contract as implemented. Controlled ablations established the contribution of each component: the multi-cycle data, the point-cloud encoder, the generative head, the human-corrected data, the Cartesian contract, and the progress and phase conditions. Online DAgger variants with recurrent regression heads, behavior cloning, and the joint-space action contract were evaluated and found insufficient for the cyclic task, and the evaluation ladder localized their failures in recovery from the student's own deviations.

Experiments show that the distilled policy completes ten consecutive threading cycles on all training-distribution nut sizes in simulation and on hardware, with mean cyclic tracking errors below $6\degree$, and transfers to unseen round and octagonal profiles, while rectangular profiles expose a limitation of the learned re-grasp strategy. These results suggest that structured behavior generation followed by unified policy distillation is an effective approach to long-horizon, contact-rich dexterous manipulation.

\section{Future Work}\label{sec:future-work}

Long-horizon stress tests reveal the first direction. The current student fails after 15 to 22 cycles because the axial drift accumulated over many cycles leaves the range covered by the demonstrations. A curriculum that progressively expands the axial variation encountered during teacher training, continuing the 20\,mm initial-height randomization toward the full usable thread, would improve the long-horizon capability of the teacher and, consequently, of the distilled student. The global progress clock should at the same time be replaced by phase inference from the observation history inside the policy, so that the student is not bound to a nominal trial duration.

The second direction is autonomous nut engagement. The task was deliberately initialized with the nut already engaged after unsuccessful attempts to learn or demonstrate engagement with reinforcement learning, teleoperation, and imitation learning; the difficulty is attributed to the narrow entry configuration that a rigid nut and the helical constraint admit. Engagement will require fine-grained fingertip-force feedback and tactile sensing; the hardware-equivalent fingertip-force and joint-torque observation channels implemented at the end of this project are the starting point, and the pickup expert's handoff and engagement goals are the open teacher problem.

The third direction concerns the distillation itself. The online flow-matching student is implemented but unqualified; collecting multi-step teacher recovery fragments from student-visited states would give it the DAgger property that the recurrent students benefited from, without violating the chunk-label constraint. The pending sweeps over the action horizon, the number of integration steps, and the temporal-ensemble decay, and the closed-loop comparison of the Adapt3R-style and point-cloud encoders, would complete the design study. Finally, the generalization study suggests extending the training-object family beyond hexagonal profiles so that the student can learn re-grasp strategies for face topologies that the hexagon does not cover.
